\IEEEraisesectionheading{\section{Introduction}\label{sec:introduction}}

In May 2022, the Terra/Luna collapse erased over \$40 billion in value within 48~hours, triggering cascading liquidations across lending protocols, DEXs, and yield aggregators that held UST-denominated assets. Five months later, the FTX insolvency exposed hidden cross-protocol exposures and wiped out another \$8 billion. These crises share a common pattern: shocks propagated through an interconnected web of tokenised liabilities---credit exposures---that no single participant fully observed. In decentralised finance, exposures are high-dimensional, time-varying, and often implicit; data are noisy and incomplete; and crises unfold within days, leaving limited time for manual assessment.

Automated \emph{agentic} pipelines offer a principled response: an agent can ingest on-chain measurements, maintain a belief over the evolving exposure network, surface early-warning signals, stress-test plausible shocks, and recommend risk-mitigating actions with auditable rationale. However, existing agentic proposals for DeFi risk largely rely on general-purpose large language models or static heuristics, and lack a dedicated forecasting backbone for graph-structured time series. Without accurate and data-efficient prediction of network dynamics, monitoring agents suffer from high false-alarm rates and delayed detection, while decision agents risk recommending actions based on brittle or poorly calibrated beliefs. In high-stakes contexts, uncertainty calibration and robustness to missingness and regime shifts are not optional---they determine whether an agent is trusted and deployable.

This paper addresses that bottleneck by formulating a forecast-grounded regulatory decision framework for DeFi exposure networks. We leverage \textsc{DeXposure-FM}, a foundation model trained on over $43.7$~million data entries from DeFi exposure networks, to forecast the evolution of credit-exposure graphs across multiple horizons. The framework then follows a \emph{forecast--measure--decide--gate} paradigm: predicted graphs are converted into risk measures and scenario evidence, decision policies turn this evidence into supervisory tickets, and data-health/confidence gates determine whether intervention-level actions are allowed. To test whether such a framework is actually useful, we introduce \textsc{DeXposure-Bench}; to show that it can be operated as an auditable system, we instantiate it as \textsc{DeXposure-Claw}.

Our contributions are:
\begin{enumerate}[leftmargin=*, itemsep=2pt]
    \item \textbf{A forecast-grounded regulatory decision framework for DeFi exposure networks} (\S\ref{sec:algorithm}). We formulate DeFi risk supervision as a four-layer pipeline: FM prediction, deterministic monitoring and scenario analysis, decision policy, and data-health/confidence gating. This separates the problem of forecasting exposure dynamics from the problem of deciding when a supervisory ticket is warranted. Five network risk metrics with rolling baselines, five standardised stress scenarios with CVaR aggregation, and a four-level playbook make the framework auditable rather than a black-box LLM decision loop. While prior agentic work in DeFi targets transaction verification~\cite{yao2025arbiter} or intent mining~\cite{mao2025knowyourintent}, and graph-based foundation models for risk have been explored only in traditional finance~\cite{li2025hgmae}, no existing framework unifies forecasting, stress evidence, decision recommendation, and fail-safe gating for DeFi exposure networks.

    \item \textbf{\textsc{DeXposure-Bench}: a benchmark for evaluating forecast-grounded risk-decision pipelines} (\S\ref{sec:bench}). To test the framework, we release a standardised harness that scores a candidate pipeline on six axes: forecasting quality, streaming anomaly detection, predictive uncertainty, what-if scenario fidelity, supervisory decision quality, and data-quality robustness. The bench operates on the full \textsc{DeXposure} corpus ($4{,}300$ protocols, $24{,}300$ tokens, $43.7$M exposure entries, $283$ weekly snapshots from $2020$ to $2025$~\cite{wu2025dexposuredatasetbenchmarksinterprotocol}), uses a regulator-aligned ground truth defined by \emph{absolute} weight loss rather than fractional change, and ships eight reference method implementations. Generic LLM-agent benchmarks (\textsc{SWE-bench}, \textsc{AgentBench}, \textsc{HELM}) do not evaluate domain-grounded supervisory decisions, and temporal-graph benchmarks (\textsc{TGB}, \textsc{OGB}) stop at structural prediction; \textsc{DeXposure-Bench} fills this benchmark gap for agentic risk pipelines.

    \item \textbf{\textsc{DeXposure-Claw}: a deployable agentic instantiation and empirical characterisation} (\S\ref{sec:exp}). We package the framework as an auditable agentic system that emits ticket-level evidence bundles and can be run under a weekly supervisory cadence. On \textsc{DeXposure-Bench}, the full FM\,+\,LLM\,+\,gate stack raises ticket F1 from $0.0076$ (persistence\,+\,rules) to $0.0233$ (\textbf{$+207\%$}). A tier-above Claude~Opus~4.8 judge with a cross-family panel (Gemini~2.5~Pro, GPT-5.5) preserves the ordering \texttt{m7}$\geq$\texttt{m6}$>$\texttt{m2}; under the Opus~4.7 decision LLM the false-intervention rate rises to $\approx 0.44$. Replacing the decision LLM with a tier-below Sonnet~4.6 \emph{strictly dominates} Opus~4.7 on both \texttt{m6} and \texttt{m7}: F1 rises further to $0.0288$ on \texttt{m7}, FIR collapses to $0.000$, and judge quality reaches $2.65/5$, at $\sim$5$\times$ lower inference cost---the deployable headline operating point.
\end{enumerate}

We will release the implementation and benchmark artifacts with the project repository.
